\chapter{The Weak Equation}
\label{chap:weak-equation}

The strong equation asks for \(R_U(\psi)=0\) pointwise. That is often too much to ask of the apparatus and too
much to ask of the field. Boundary traces may exist only weakly. Derivatives may live in distributional form.
The domain may be known only through finite elements. The measurement may see the residual only after it has
been integrated against test functions. The weak equation is the honest equation at that floor:
\[
  \langle R_U(\psi), \eta\rangle_{H^s} = 0
  \qquad
  \hbox{for all admissible }\eta .
\]

\section{Strong Enough to Compute}

Weak does not mean vague. It means the residual is read through tests. If the tests are rich enough, the weak
statement is as strong as the experiment can honestly support. In Galerkin form, the admissible tests are a
finite basis \(\{\phi_i\}_{i=1}^{N}\). The state is approximated in the same or a companion finite space,
\[
  \psi_N = \sum_{j=1}^{N} c_j \phi_j,
\]
and the weak equations become
\[
  \langle R_U(\psi_N), \phi_i\rangle_{H^s}=0,
  \qquad i=1,\ldots,N.
\]
That is a linear or nonlinear algebraic system depending on \(R_U\). In the linearized skeleton here, it is an
exact finite system. The serious part is not the size of the matrix. The serious part is the contract: every
test row is a weak equation, and satisfying all rows is the finite statement that the first variation vanishes
on the chosen space.

\section{The Exact Galerkin Predicate}

The Lean file names this finite closure:

\begin{quote}\small
\begin{verbatim}
def ExactGalerkin
    (U : UniversalTensor) (state : Vec) : Prop :=
  sobolevResidual U state =
    zeroLike (sobolevResidual U state)
\end{verbatim}
\end{quote}

This predicate is intentionally stricter than a tolerance. It says the Sobolev residual vector is exactly zero
in the finite arithmetic of the chosen rows. In a floating-point production code one may use residual norms,
condition estimates, and stopping criteria. But the mathematical skeleton should know what exact closure means
before numerical tolerance enters. The toy computation uses integer coefficients so the zero is literal.

The proof that exact Galerkin closure implies the weak equation is the core clasp:

\begin{quote}\small
\begin{verbatim}
theorem exactGalerkin_is_weakLeastActivity
    (U : UniversalTensor) (state : Vec)
    (h : ExactGalerkin U state) :
    WeakLeastActivity U state := by
  intro test
  unfold ExactGalerkin at h
  unfold firstFrechetVariation
  rw [h]
  exact dot_zero_left_any (sobolevResidual U state) test
\end{verbatim}
\end{quote}

The proof is short because the definitions have been set in the right order. If the Sobolev residual is the zero
vector, then its dot product with every finite test vector is zero. Therefore the first Fréchet variation
vanishes on every finite test vector. Therefore the state is a weak least-activity solution at the chosen
Galerkin floor.

\section{What the Theorem Says}

The theorem does not say that every Dirac equation has been solved. It does not say that every Sobolev space has
been constructed. It does not say that regularity, compactness, ellipticity, or convergence have been obtained
for free. Those belong to a full analysis text. The theorem says something narrower and load-bearing: after the
Dirac rows, tensor rows, Sobolev reader, load, state, and test space have been fixed, exact finite Galerkin
closure is enough to discharge the weak first-variation equation on that finite space.

That is exactly the floor Book II needs. Experimentation is not allowed to gesture at the continuum and then
forget the rows. It must say what it computed. The weak theorem says what was computed: a finite residual,
tested in a Sobolev reader, vanishing against every finite test vector.

\section{Dirac Inside the Weak Gate}

The Dirac equation enters through the residual rows, not as a slogan. In a real discretization the rows contain
the Clifford matrices, the gauge connection, mass coupling, boundary conditions, and any tensor background. The
finite vector \(\psi_N\) stores the coefficients of the spinor in the chosen basis. The weak equation asks that
each test spinor see no Sobolev residual:
\[
  \left\langle
    \mathcal D_A\psi_N + \mathcal T_U\psi_N - f,\,
    \phi_i
  \right\rangle_{H^s}
  = 0 .
\]
This is where the Dirac operator becomes computational rather than ornamental. It is the first-order term whose
rows are assembled, tested, and closed.
